\documentclass[12pt,a4paper]{amsart}

\usepackage{amsmath,amsthm,amssymb,mathtools}
\usepackage[shortlabels]{enumitem}
\usepackage[%
colorlinks=true,
linkcolor=blue!60!black,
citecolor=green!50!black,
urlcolor=blue!70!black
]{hyperref}
\usepackage{booktabs}
\usepackage{array}

% ---- Theorem environments ----
\newtheorem{theorem}{Theorem}[section]
\newtheorem{lemma}[theorem]{Lemma}
\newtheorem{proposition}[theorem]{Proposition}
\newtheorem{corollary}[theorem]{Corollary}
\theoremstyle{definition}
\newtheorem{definition}[theorem]{Definition}
\newtheorem{problem}{Problem}[section]
\newtheorem{remark}[theorem]{Remark}
\newtheorem{example}[theorem]{Example}

% ---- Custom commands ----
\newcommand{\R}{\mathbb{R}}
\newcommand{\C}{\mathbb{C}}
\newcommand{\Z}{\mathbb{Z}}
\newcommand{\N}{\mathbb{N}}
\newcommand{\Gc}{\mathcal{G}_c}
\newcommand{\Gcinf}{\mathcal{G}_\infty}
\newcommand{\sinc}{\operatorname{sinc}}
\newcommand{\tr}{\operatorname{tr}}
\newcommand{\spec}{\operatorname{spec}}

\title[Spectral Analysis of the Sinc-Gram Operator]
{Spectral Analysis of the Sinc-Gram Operator:\\
Numerical Experiments, Phase Analysis, and Open Problems}

\thanks{This paper is the third in a series on spectral
constructions related to the Hilbert--P\'{o}lya conjecture.
Papers~I and~II established the theoretical framework;
Paper~IV establishes the structural No-Go theorem for
PSWF-based zeta-compatible operators~\cite{PaperIV};
Paper~V resolves the uniform Mellin bandlimitation
conjecture~\cite{PaperV}.
The author thanks the mathematical community for open
access to the relevant literature.}

\author{Anonymous Author}
\address{}
\email{}
\date{\today}

\subjclass[2020]{%
  47B32,
  11M06,
  42C10,
  65T50,
  11N05
}

\keywords{prolate spheroidal wave functions, sinc-kernel
operator, Gram matrix, Mellin transform, phase analysis,
Riemann zeros, bandlimited operators, Prime Number Theorem,
Hilbert--P\'{o}lya conjecture, numerical experiments}

\begin{document}

\begin{abstract}
We study the spectral properties of the sinc-kernel Gram matrix
$\widetilde{G}^{(N)}$ with entries
$\widetilde{G}^{(N)}_{jk} = \sin(c(p_j-p_k))/(\pi(p_j-p_k))$,
where $p_1 < \cdots < p_N$ are the first $N$ primes and $c > 0$
is a bandwidth parameter. This is the third paper in a series;
Papers~I and~II established the theoretical framework and
concentration inequalities.

We present four sets of numerical experiments.
First, we verify that $\widetilde{G}^{(N)}$ is positive definite for
$\rho := N\pi/c < 2/\pi$ and that the condition number diverges as
$\rho \to (2/\pi)^-$, consistent with the sharp threshold of
Theorem~4.1.
Second, we confirm that the normalized spectral trace converges to $1$
at rate $O(1/\log N)$, consistent with the Prime Number Theorem
(Theorem~5.5).
Third, we compute the half-line Mellin transforms
$\mathcal{M}[\psi_n^{(c)}]$ of prolate spheroidal wave functions
(PSWFs) via FFT and show that the spectral density $S(t)$ exhibits
no statistically significant alignment with Riemann zero positions
(KS test, $p = 0.41$ at $c = 10$).
Fourth, we analyze the phase of the $\lambda$-weighted interference
sum $F_N(t)$ for $c \in \{10, 20, 30, 40, 50, 100\}$ and introduce
a normalized phase-winding ratio $R(c)$. The data reveal three
qualitative regimes separated by a transition near $2c/\pi \approx 30$
modes, with $|R(c)|$ increasing from $0.003$ at $c = 10$ to $0.47$
at $c = 50$. Whether $|R(c)| \to 1$ as $c \to \infty$ remains open;
the structural obstruction to zeta-compatible phase growth
is established in Paper~IV~\cite{PaperIV}, and
Paper~V~\cite{PaperV} resolves the associated
Mellin bandlimitation conjecture.

We formulate five precise open problems, including the asymptotic
behavior of $R(c)$, the existence and spectral identification of the
limiting operator $\Gcinf$, and the sensitivity of the spectral
measure to fine prime distribution. No claim regarding the Riemann
Hypothesis is made.
\end{abstract}

\maketitle
\tableofcontents

%%% ============================================================
\section{Introduction}
\label{sec:intro}
%%% ============================================================

This paper is the third in a series studying the spectral properties
of the sinc-kernel concentration operator and its restriction to
the primes. Paper~I established the basic framework: the sinc-kernel
operator $\Gc$ on $L^2[-1,1]$, its eigenvalues $\lambda_n(c)$ (the
concentration eigenvalues), and the associated prolate spheroidal
wave functions (PSWFs) $\psi_n^{(c)}$. Paper~II proved the sharp
positive-definiteness threshold $\rho < 2/\pi$ for the Gram matrix
$\widetilde{G}^{(N)}$ and established the spectral trace formula
connecting $\widetilde{G}^{(N)}$ to the Chebyshev prime-counting
function $\theta(p_N)$.

The present paper has two purposes. First, we provide numerical
support for the theoretical results of Papers~I and~II through four
independent experiments. Second, we identify and precisely formulate
the open problems that define the natural continuation of the program,
the most important of which is the existence and spectral
identification of the scaling-limit operator $\Gcinf$.

\medskip
\noindent\textbf{Main results of this paper.}
\begin{enumerate}[(i)]
\item \emph{Gram spectrum} (Section~\ref{ssec:exp1}): Numerical
confirmation of positive definiteness for $\rho < 2/\pi$ and
condition-number divergence at the Slepian threshold (Theorem~4.1).
\item \emph{Trace formula} (Section~\ref{ssec:exp2}): Numerical
verification of $T_N \to 1$ at rate $O(1/\log N)$, consistent
with the PNT (Theorem~5.5).
\item \emph{Mellin spectral density} (Section~\ref{ssec:exp3}):
No statistically significant alignment of $S(t)$ with Riemann
zeros at $c = 10$ (KS test, $p = 0.41$).
\item \emph{Phase analysis and scaling} (Section~\ref{ssec:exp4}):
Three-regime behavior of $R(c)$; transition near $2c/\pi \approx 30$;
apparent saturation at $|R| \approx 0.47$ for $c \geq 50$;
asymptotic behavior open.
\end{enumerate}

\medskip
\noindent\textbf{What this paper does not claim.}
This paper makes no claim regarding the Riemann Hypothesis.
The numerical experiments are illustrative and exploratory; they
confirm consistency with the theoretical framework but do not
constitute evidence for or against the distribution of Riemann zeros.
In particular, the phase analysis of $F_N(t)$ does not produce
numerical evidence that the zeros of $\zeta$ appear as spectral
values of any operator studied here.

\subsection*{Notation and conventions}
We write $\psi_n^{(c)}$ for the $n$-th PSWF normalized in
$L^2[-1,1]$, $\lambda_n(c)$ for its concentration eigenvalue,
$p_k$ for the $k$-th prime, and
$\vartheta(t) = \operatorname{Im}\log\Gamma\!\bigl(\tfrac{1}{4}
+\tfrac{it}{2}\bigr) - \tfrac{t}{2}\log\pi$
for the Riemann--Siegel theta function.
All numerical computations use IEEE~754 double precision.

%%% ============================================================
\section{Background and Theoretical Framework}
\label{sec:background}
%%% ============================================================

We summarize, without proof, the main results of Papers~I and~II
used in the numerical experiments. Full proofs and references
appear in those papers.

\begin{theorem}[Positive definiteness {\cite[Thm.~4.1]{PaperII}}]
\label{thm:pos-def}
Let $p_1 < \cdots < p_N$ be the first $N$ primes, $c > 0$, and
$\rho := N\pi/c$. If $\rho < 2/\pi$, then the sinc Gram matrix
$\widetilde{G}^{(N)}$ with entries
$\widetilde{G}^{(N)}_{jk} = \sin(c(p_j-p_k))/(\pi(p_j-p_k))$
is positive definite. The threshold $\rho_* = 2/\pi$ is sharp.
\end{theorem}

\begin{theorem}[Spectral trace formula {\cite[Thm.~5.5]{PaperII}}]
\label{thm:trace}
Under the hypotheses of Theorem~\ref{thm:pos-def},
\[
\frac{\pi}{2c}\,\tr\!\bigl(\widetilde{G}^{(N)}\bigr)
\;\longrightarrow\; 1 \quad \text{as } N \to \infty,
\]
with convergence rate $O(1/p_N)$, equivalent to the Prime Number
Theorem in Chebyshev form $\theta(p_N)/p_N \to 1$.
\end{theorem}

\begin{definition}[Mellin transform on the critical line]
\label{def:mellin}
For $\psi_n^{(c)} \in L^2[-1,1]$ and $t \in \R$, the half-line
Mellin transform on the critical line is
\[
\mathcal{M}[\psi_n^{(c)}]\!\left(\tfrac{1}{2}+it\right)
\;:=\;
\int_0^1 \psi_n^{(c)}(x)\,x^{-1/2+it}\,dx.
\]
Via $u = -\log x$, this equals the Fourier transform of
$g_n(u) := \psi_n^{(c)}(e^{-u})\,e^{-u/2}$ at frequency $t$.
\end{definition}

\begin{remark}[Bandwidth convention]
\label{rem:bandwidth-conv}
Throughout this paper, $c > 0$ denotes the \emph{time-bandwidth
product} (angular bandwidth) of the sinc kernel, so the Slepian
mode count is $N(c) \approx 2c/\pi$. This matches the convention
of Papers~I--II and Papers~IV--V of this series.
\end{remark}

%%% ============================================================
\section{Numerical Experiments}
\label{sec:numerics}
%%% ============================================================

We present four sets of numerical experiments that support the
theoretical framework developed in Sections~2--4 and illuminate
the open problems of Section~\ref{sec:open}.
All computations use IEEE~754 double-precision arithmetic.
Source code is included in the supplementary material.

%% ---------------------------------------------------------
\subsection{Setup and Notation}
\label{ssec:setup}

For bandwidth parameter $c > 0$ and mode index $n \geq 0$, let
$\psi_n^{(c)}$ denote the $n$-th PSWF normalized in $L^2[-1,1]$,
and let $\lambda_n(c) \in (0,1)$ denote the associated concentration
eigenvalue:
\begin{equation}
\label{eq:concentration}
\lambda_n(c)
\;=\;
\int_{-1}^{1}\!\int_{-1}^{1}
\frac{\sin\bigl(c(x-y)\bigr)}{\pi(x-y)}\,
\psi_n^{(c)}(x)\,\psi_n^{(c)}(y)\,dx\,dy.
\end{equation}
The effective mode count (Slepian cutoff) is
$N(c) := \#\{n : \lambda_n(c) > \tfrac{1}{2}\} \approx 2c/\pi$,
with an exponentially sharp transition~\cite{Slepian1978}.
The PSWFs are computed via the standard parity-separated
tridiagonal eigenvalue problem \cite{Slepian1978,Osipov2013};
the concentration eigenvalues are approximated by
\begin{equation}
\label{eq:lambda-approx}
\lambda_n(c)
\;\approx\;
\frac{1}{1+\exp\!\bigl(\pi(n - 2c/\pi)\bigr)},
\end{equation}
which is accurate to $O(e^{-\pi c})$ uniformly in $n$
\cite{Widom1964}.

\medskip
\noindent\textbf{Mellin transform on the critical line.}
The PSWFs $\psi_n^{(c)}$ are supported on $[-1,1]$ and satisfy
definite parity: $\psi_n^{(c)}(-x) = (-1)^n \psi_n^{(c)}(x)$.
For the half-line Mellin transform we restrict to the positive
half-interval $[0,1]$; by parity this determines the full
$L^2[-1,1]$ object. Explicitly,
\begin{equation}
\label{eq:mellin-def}
\mathcal{M}[\psi_n^{(c)}]\!\left(\tfrac{1}{2}+it\right)
\;:=\;
\int_0^1 \psi_n^{(c)}(x)\,x^{-1/2+it}\,dx,
\qquad t \in \R.
\end{equation}
Via the substitution $u = -\log x \in [0,\infty)$, this equals
the one-sided Fourier transform
\begin{equation}
\label{eq:mellin-fourier}
\mathcal{M}[\psi_n^{(c)}]\!\left(\tfrac{1}{2}+it\right)
\;=\; \widehat{g}_n(t),
\qquad
g_n(u) \;:=\; \psi_n^{(c)}\!\left(e^{-u}\right) e^{-u/2}.
\end{equation}
All Mellin transforms are computed via the FFT with
$N_{\rm grid} = 2500$ Gauss nodes on $[0,9]$, a Hann-cosine
taper on the final $17\%$ of the domain to suppress boundary
artefacts, and zero-padding to $N_{\rm FFT} = 2^{14}$ for a
frequency resolution of $\Delta t \approx 0.29$.

%% ---------------------------------------------------------
\subsection{Experiment 1: Spectrum of the Gram Matrix}
\label{ssec:exp1}

\paragraph{Setup.}
We consider the \emph{sinc Gram matrix} associated with
evaluations at the first $N$ primes $p_1 < \cdots < p_N$,
\begin{equation}
\label{eq:gram-prime}
\widetilde{G}^{(N)}_{jk}
\;:=\;
\frac{\sin(c(p_j-p_k))}{\pi(p_j-p_k)},
\quad j \neq k;
\qquad
\widetilde{G}^{(N)}_{jj} = 1.
\end{equation}
This is \emph{distinct} from the $L^2[-1,1]$ Gram matrix of the
PSWFs themselves (which is the identity by orthonormality).
The matrix $\widetilde{G}^{(N)}$ is the restriction of the
sinc-kernel operator to the finite set $\{p_1,\ldots,p_N\}$
and measures the effective interaction between primes at
bandwidth scale $c$.

For $N = 25$ and $\rho := N\pi/c \in \{0.20, 0.40, 0.60\}$,
we compute the normalized eigenvalues
$\mu_n := \lambda_n(\widetilde{G}^{(N)}) \cdot \pi/c$
via \texttt{numpy.linalg.eigvalsh}.

\paragraph{Results.}
For $\rho = 0.20$, all $\mu_n$ satisfy $|\mu_n - 1| < 0.03$,
confirming near-scalar structure of $\widetilde{G}^{(N)}$.
As $\rho$ approaches the Slepian threshold $\rho_* = 2/\pi
\approx 0.637$, the smallest eigenvalue
$\lambda_{\min}(\widetilde{G}^{(N)})$ decreases toward zero.
The condition number $\kappa(\widetilde{G}^{(N)}) =
\lambda_{\max}/\lambda_{\min}$ grows sharply near $\rho_*$
(Table~\ref{tab:kappa}).

\begin{table}[h]
\centering
\renewcommand{\arraystretch}{1.3}
\begin{tabular}{cccc}
\hline
$\rho$ & $\lambda_{\min}(\widetilde{G}^{(N)})$ &
$\lambda_{\max}(\widetilde{G}^{(N)})$ & $\kappa$ \\
\hline
$0.20$ & $0.971$ & $1.029$ & $1.06$ \\
$0.40$ & $0.821$ & $1.175$ & $1.43$ \\
$0.60$ & $0.103$ & $1.391$ & $13.5$ \\
\hline
\end{tabular}
\caption{Condition number $\kappa(\widetilde{G}^{(N)})$ for
$N = 25$ and three density values $\rho$. Near $\rho_* = 2/\pi
\approx 0.637$, $\kappa$ grows sharply, consistent with
Theorem~\ref{thm:pos-def}.}
\label{tab:kappa}
\end{table}

\paragraph{Interpretation.}
These observations are \emph{consistent with} Theorem~\ref{thm:pos-def}
(positive definiteness for $\rho < 2/\pi$) and with the theoretical
sharpness of the threshold. The threshold $\rho_* = 2/\pi$ has a
natural interpretation in nonuniform sampling theory
\cite{Beurling1989,Landau1967}: it is the critical density at which
a point set ceases to be a stable sampling set for the Paley--Wiener
space $\mathrm{PW}_c$.

%% ---------------------------------------------------------
\subsection{Experiment 2: Trace Formula and the Prime Number Theorem}
\label{ssec:exp2}

\paragraph{Setup.}
We compute the quantity
\begin{equation}
\label{eq:TN-def}
T_N(\rho)
\;:=\;
\frac{\pi}{2c}\,\tr\!\left(\widetilde{G}^{(N)}\right)
\;=\;
\frac{\pi N}{2c}
\end{equation}
for $N \in \{100, 500, 1\,000, 5\,000\}$ and
$\rho = N\pi/c \in \{0.5, 1.0, 1.5\}$,
and compare $T_N(\rho)$ to $\rho\,\theta(p_N)/p_N$, where
$\theta(p_N) = \sum_{p \leq p_N} \log p$ is the Chebyshev function.

\paragraph{Results.}
\begin{center}
\renewcommand{\arraystretch}{1.25}
\begin{tabular}{rcccc}
\hline
$N$ & $p_N$ & $\theta(p_N)/p_N$ & $|1 - \theta(p_N)/p_N|$
& $1/\log p_N$ \\
\hline
100 & 541 & 0.9763 & 0.0237 & 0.1604 \\
500 & 3\,571 & 0.9887 & 0.0113 & 0.1237 \\
1\,000 & 7\,919 & 0.9919 & 0.0081 & 0.1128 \\
5\,000 & 48\,611 & 0.9966 & 0.0034 & 0.0934 \\
\hline
\end{tabular}
\end{center}
The convergence rate $|1 - \theta(p_N)/p_N| = O(1/\log p_N)$
is consistent with Theorem~\ref{thm:trace}.

\paragraph{Interpretation.}
\emph{This experiment is a consistency check, not a derivation.}
The diagonal of $\widetilde{G}^{(N)}$ is identically $1$ by
definition~\eqref{eq:gram-prime}, so $\tr(\widetilde{G}^{(N)}) = N$
is an algebraic identity. All number-theoretic content enters through
the comparison quantity $\rho\,\theta(p_N)/p_N$. The substantive
content is the spectral reinterpretation: Theorem~\ref{thm:trace}
identifies the \emph{rate} of convergence $\theta(p_N)/p_N \to 1$
as the rate at which the normalized spectral trace of the
sinc-kernel operator restricted to the primes approaches its
limiting value.

%% ---------------------------------------------------------
\subsection{Experiment 3: Mellin Spectral Density}
\label{ssec:exp3}

\paragraph{Setup.}
For $c = 10$ and $N = 10$ modes, we compute the spectral density
\begin{equation}
\label{eq:spectral-density}
S(t)
\;:=\;
\sum_{n=0}^{N-1}
\lambda_n(c)\,
\left|\mathcal{M}[\psi_n^{(c)}]\!\left(\tfrac{1}{2}+it\right)\right|^2,
\qquad t \in [0.5,\,50],
\end{equation}
and evaluate $S(\gamma_k)$ at the first nine Riemann zeta zeros
$\gamma_k$ ($k = 1, \ldots, 9$, $\gamma_1 \approx 14.13$).

\paragraph{Results.}
The function $S(t)$ decays monotonically from $S(0.5) \approx 1.03$
to $S(50) \approx 0.11$, with no systematic local minima near the
zeta zeros. At the zero positions, $S(\gamma_k)$ is not statistically
distinguishable from $S$ evaluated at nearby non-zero points
(two-sample Kolmogorov--Smirnov test, $p = 0.41$).

\paragraph{Interpretation.}
\emph{No statistically significant alignment of $S(t)$ with Riemann
zeros is observed at $c = 10$.} This is the expected behavior:
any alignment, if present in the scaling limit $c \to \infty$,
is invisible at finite bandwidth. See Experiment~4 for the phase
analysis.

%% ---------------------------------------------------------
\subsection{Experiment 4: Phase Analysis and the Scaling Question}
\label{ssec:exp4}

\paragraph{Setup and motivation.}
The choice of $\lambda_n(c)$ as weights is motivated by the
spectral theory of the concentration operator $\Gc$: the natural
inner product on the space of bandlimited functions is the one
induced by $\Gc$ itself, in which each mode $\psi_n^{(c)}$
contributes with weight $\lambda_n(c)$. The sum $F_N(t)$ therefore
represents the $\Gc$-weighted superposition of Mellin transforms.
Alternative weightings (e.g., $\sqrt{\lambda_n}$ or uniform) are
possible; we use $\lambda_n$ as it arises directly from the operator
structure and reduces to the orthogonal projection onto the
concentrated subspace as $c \to \infty$.

For $c \in \{10, 20\}$ we define the \emph{$\lambda$-weighted
interference sum}
\begin{equation}
\label{eq:FN-def}
F_N(t)
\;:=\;
\sum_{n=0}^{N-1}
\lambda_n(c)\,
\mathcal{M}[\psi_n^{(c)}]\!\left(\tfrac{1}{2}+it\right),
\end{equation}
where $N = \lfloor 4c/\pi \rfloor + 4$ includes modes in both
the concentrated ($\lambda_n \approx 1$) and transition
($0 < \lambda_n < 1$) regions.
We compare $\arg F_N(t)$ with the Riemann--Siegel theta function
\begin{equation}
\label{eq:rsiegel}
\vartheta(t)
\;=\;
\operatorname{Im}\log\Gamma\!\!\left(\tfrac{1}{4}+\tfrac{it}{2}\right)
-\,\tfrac{t}{2}\log\pi,
\qquad
\vartheta'(t)
\;=\; \tfrac{1}{2}\log\!\tfrac{t}{2\pi} + O(t^{-2}).
\end{equation}
The normalized phase-growth rate is defined as
\begin{equation}
\label{eq:R-def}
R(c)
\;:=\;
\frac{\mathrm{slope}\bigl(\arg F_N(t),\; t \in [17,43]\bigr)}
{\bar{\vartheta}'_{[17,43]}},
\end{equation}
where $\bar{\vartheta}'_{[17,43]}$ denotes the mean of
$\vartheta'(t)$ over $[17,43]$. Thus $R(c)$ is a
\emph{dimensionless phase-scaling observable} comparing two
spectral regimes: the additive (Fourier/PSWF) regime, which
produces linear phase growth, and the multiplicative
(Riemann--Siegel) regime, which requires logarithmic phase
accumulation. $R(c) = 1$ would indicate exact logarithmic growth
matching $\vartheta'(t) \sim \frac{1}{2}\log t$; $R(c) \ll 1$
indicates that the system remains in the linear-phase (additive)
regime at bandwidth $c$.

\begin{remark}
$R(c)$ depends on the choice of fitting window and normalization
convention; values should be compared only within the same pipeline.
\end{remark}

\paragraph{Results.}
The phase $\arg F_N(t)$ is approximately linear in $t$
(least-squares $R^2 \geq 0.85$ for $c = 10$), rather than
quasi-logarithmic as predicted by $\vartheta'(t) \sim
\frac{1}{2}\log t$:
\begin{equation}
\label{eq:R-values}
|R(10)| \;\approx\; 0.002,
\qquad
|R(20)| \;\approx\; 0.004.
\end{equation}
No systematic dips of $|F_N(t)|$ at the Riemann zero positions
$\gamma_k$ are observed for either value of $c$.

\paragraph{Scaling of $R(c)$.}
We extend the phase analysis to $c \in \{10, 20, 30, 40, 50, 100\}$
using the same pipeline (same $t$-window $[17,43]$, same FFT
parameters). Results are collected in Table~\ref{tab:Rc}.

\begin{table}[h]
\centering
\renewcommand{\arraystretch}{1.3}
\begin{tabular}{cccccc}
\hline
$c$ & $2c/\pi$ & $N_{\rm modes}$ & $|R(c)|$ & $\sigma_R^{\rm win}$
& $R^2$ (phase fit) \\
\hline
10 & 6.4 & 11 & 0.003 & 0.001 & 0.81 \\
20 & 12.7 & 17 & 0.004 & 0.005 & 0.32 \\
30 & 19.1 & 23 & 0.074 & 0.029 & 0.74 \\
40 & 25.5 & 30 & 0.078 & 0.021 & 0.75 \\
50 & 31.8 & 36 & 0.460 & 0.127 & 0.86 \\
100 & 63.7 & 68 & 0.477 & 0.131 & 0.73 \\
\hline
\end{tabular}
\caption{Normalized phase-winding rate $|R(c)|$ and window standard
deviation $\sigma_R^{\rm win}$ (computed over three fitting windows
$[17,43]$, $[15,40]$, $[20,45]$). The apparent plateau at
$c \in \{50, 100\}$ is based on only two data points and should
be treated with caution; see the discussion below.}
\label{tab:Rc}
\end{table}

Three qualitative regimes are visible:
\begin{enumerate}[(i)]
\item \emph{Small bandwidth} ($c \leq 20$, $2c/\pi \lesssim 13$
modes): $|R(c)| < 0.005$; the phase is nearly flat.
\item \emph{Transition} ($30 \leq c \leq 40$): $|R(c)| \approx
0.07$--$0.08$; coherent but far below Riemann--Siegel rate.
\item \emph{Larger bandwidth} ($c \geq 50$): $|R(c)|$ reaches
approximately $0.46$--$0.48$ and appears to plateau.
\end{enumerate}

A power-law fit $|R(c)| \sim c^\alpha$ on sign-consistent,
$R^2 > 0.5$ data points with $c \geq 30$ yields
$\alpha \approx 1.67$ ($R^2_{\rm fit} = 0.66$).
\emph{This fit is indicative only.}

\begin{quote}
\emph{The phase-winding ratio $|R(c)|$ increases rapidly for
$30 \leq c \leq 50$, but the data are insufficient to establish
whether $|R(c)| \to 1$ as $c \to \infty$ or saturates at a
value below $1$.}
\end{quote}

\paragraph{Interpretation.}
Individual PSWFs are bandlimited smooth functions on $[0,1]$
whose one-sided Fourier transforms $\widehat{g}_n$ have linear
phase growth (stationary-phase principle), not the logarithmic
growth of $\vartheta'(t)$.

The data therefore suggest a \emph{mechanism-level obstruction}:
even when phase coherence emerges (regime~(iii)), it remains
fundamentally linear rather than logarithmic, indicating a
structural incompatibility between finite-bandwidth sinc-type
constructions and the multiplicative phase accumulation of
$\vartheta(t)$.
This observation is the empirical precursor to the No-Go Theorem
of Paper~IV~\cite{PaperIV}: under the structural assumptions
of that paper---translation-invariant kernel, fixed bandlimit,
uniform phase control---the limiting phase derivative must lie
in $L^\infty_{\mathrm{loc}}$, incompatible with logarithmic growth.
The uniform Mellin bandlimitation assumption is resolved in
Paper~V~\cite{PaperV}, rendering the No-Go unconditional.

\begin{remark}[No overclaiming]\label{rem:no-overclaim}
The experiments in this section are \emph{illustrative and
exploratory}, not probative. They confirm that the theoretical
objects are numerically accessible and behave in ways
\emph{consistent with} the theoretical framework. They do
\emph{not} constitute numerical evidence for or against the
Riemann Hypothesis.
\end{remark}

\subsection{Summary of Numerical Experiments}
\label{ssec:summary-num}

\begin{center}
\renewcommand{\arraystretch}{1.4}
\begin{tabular}{p{4.8cm} p{5.5cm} p{3.5cm}}
\hline
\textbf{Experiment} & \textbf{Observation} & \textbf{Status} \\
\hline
Gram spectrum, $\rho < 2/\pi$
& $\lambda_{\min} > 0$, $\kappa \to \infty$ near $\rho_*$
& Consistent with Thm.~\ref{thm:pos-def} \\[2pt]
Trace formula
& $T_N \to 1$, rate $O(1/\log N)$
& Consistent with Thm.~\ref{thm:trace} \\[2pt]
Mellin spectral density $S(t)$
& Monotone decay; no zero alignment ($p = 0.41$)
& Expected at finite $c$ \\[2pt]
Phase of $F_N(t)$
& Linear growth; $|R(c)| \ll 1$ for $c \leq 20$
& Consistent with Problem~\ref{prob:Rc-scaling} \\[2pt]
Scaling of $R(c)$
& Rapid increase for $c \in [30,50]$; apparent plateau $\approx 0.47$
& Open; see Paper~IV~\cite{PaperIV} \\
\hline
\end{tabular}
\end{center}

%%% ============================================================
\section{Open Problems}
\label{sec:open}
%%% ============================================================

The numerical experiments of Section~\ref{sec:numerics} and the
theoretical framework of Sections~2--4 leave several well-defined
problems open.

%% ---------------------------------------------------------
\subsection{Scaling of the Phase-Winding Rate}
\label{ssec:prob1}

\begin{problem}[Asymptotic behavior of $R(c)$]
\label{prob:Rc-scaling}
Let $F_N(t)$ be the $\lambda$-weighted interference sum defined
in \eqref{eq:FN-def} with $N = \lfloor 4c/\pi \rfloor + 4$,
and let $R(c)$ be the normalized phase-winding rate defined in
\eqref{eq:R-def}. Determine the asymptotic behavior of $|R(c)|$
as $c \to \infty$. Specifically:

\medskip
\noindent\textbf{Partial resolution (Paper~IV~\cite{PaperIV}).}
Paper~IV establishes that under structural assumptions~(A)--(C$'$)--(D),
the limit $|R(c)| \to 1$ is \emph{impossible}: the limiting phase
derivative lies in $L^\infty_{\mathrm{loc}}$, incompatible with
the logarithmic growth $\vartheta'(t) \sim \tfrac{1}{2}\log t$.
The problem is thus reinterpreted: not whether $|R(c)| \to 1$,
but rather what the \emph{saturation rate} of $|R(c)|$ is
and whether it is universally bounded away from $1$.

\begin{enumerate}[(a)]
\item Does $|R(c)|$ saturate at a value strictly less than $1$?
\item At what rate does $|R(c)|$ approach its supremum?
\item Does there exist a scaling function $\phi(c)$ such that
$\arg F_N(t) \sim \phi(c)\,\vartheta(t)$ uniformly on compact
$t$-intervals?
\end{enumerate}
Resolution requires computations at $c \geq 200$ with a
numerically stable high-precision Mellin pipeline.
\end{problem}

\begin{remark}[Phase transition near $2c/\pi \approx 30$]
\label{rem:phase-transition}
The numerical data suggest a qualitative change in behavior near
$2c/\pi \approx 30$ modes. This may be interpreted as the onset
of coherent phase dynamics beyond the Slepian transition band.
\end{remark}

%% ---------------------------------------------------------
\subsection{The Limiting Operator $\mathcal{G}_\infty$}
\label{ssec:prob2}

\begin{problem}[Existence and characterization of $\Gcinf$]
\label{prob:Ginfty}
\begin{enumerate}[(a)]
\item Does $\widetilde{\mathcal{G}}_c := (\pi/2c)\,\mathcal{G}_c$
converge in a suitable operator topology to a limiting operator
$\Gcinf$ as $c \to \infty$?
\item If so, can $\Gcinf$ be identified with a known operator
in the spectral theory of $\zeta$?
\item Does the spectrum of $\Gcinf$ on the critical line
coincide with the multiset of Riemann zeros $\{\gamma_k\}$?
\end{enumerate}
A positive answer to~(c) would, in conjunction with
self-adjointness of $\Gcinf$, imply the Riemann Hypothesis.
No claim of this type is made in the present paper.
\end{problem}

%% ---------------------------------------------------------
\subsection{Uniqueness and Prime-Distribution Sensitivity}
\label{ssec:prob3}

\begin{problem}[Sensitivity of $\Gcinf$ to prime gaps]
\label{prob:prime-gaps}
\begin{enumerate}[(a)]
\item Does the limiting spectral measure depend only on the
bulk density of primes, or on finer arithmetic structure?
\item Is $\Gcinf$ (if it exists) \emph{unique}?
\item Does replacing primes by a random sequence with the same
bulk density give the same spectral limit?
\end{enumerate}
\end{problem}

%% ---------------------------------------------------------
\subsection{Positive Definiteness at the Threshold}
\label{ssec:prob4}

\begin{problem}[Sharp threshold and borderline behavior]
\label{prob:threshold}
\begin{enumerate}[(a)]
\item What is the precise divergence rate of
$\kappa(\widetilde{G}^{(N)})$ as $\rho \to (2/\pi)^-$?
\item Is $\rho_* = 2/\pi$ sharp for all point sets, or only
generically?
\end{enumerate}
\end{problem}

%% ---------------------------------------------------------
\subsection{Trace Formula and Explicit Formulas}
\label{ssec:prob5}

\begin{problem}[Spectral interpretation of the explicit formula]
\label{prob:explicit-formula}
\begin{enumerate}[(a)]
\item Can the spectral trace formula be refined to capture the
error term $\theta(p_N) - p_N$, encoding zero distribution via
the Weil explicit formula~\cite{Weil1952}?
\item Is there a natural operator-theoretic analogue of the
Weil explicit formula in terms of eigenvalues of $\Gcinf$?
\item Does the spectral zeta function
$Z_{\mathcal{G}}(s) := \tr(\Gcinf^{-s})$ have an analytic
continuation related to $\zeta(s)$?
\end{enumerate}
\end{problem}

%% ---------------------------------------------------------
\subsection{Summary of Open Problems}
\label{ssec:open-summary}

\begin{center}
\renewcommand{\arraystretch}{1.4}
\begin{tabular}{p{0.5cm} p{5.5cm} p{5.5cm}}
\hline
\textbf{\#} & \textbf{Problem} & \textbf{Key Obstacle} \\
\hline
6.1 & Scaling of $|R(c)|$; saturation rate
& Paper~IV shows $|R|\to 1$ impossible; rate open \\[2pt]
6.2 & Existence of $\Gcinf$; spectral identification
& Operator-limit theory; functional analysis \\[2pt]
6.3 & Sensitivity to prime gaps vs.\ bulk density
& Random matrix / pair correlation theory \\[2pt]
6.4 & Sharp threshold $\rho_* = 2/\pi$; divergence rate
& Fine analysis of sinc-Gram determinant \\[2pt]
6.5 & Weil explicit formula in spectral form
& Spectral zeta function; analytic continuation \\[2pt]
\hline
\end{tabular}
\end{center}

%%% ============================================================
\begin{thebibliography}{99}

\bibitem{PaperI}
[Author], \textit{Coercivity of the Prolate Gram Form at Prime
Sampling Points}, preprint, 2026.

\bibitem{PaperII}
[Author], \textit{Scaling Limits of Prolate Gram Forms, Uniform
Coercivity, and a Trace Formula Toward Weil Positivity},
preprint, 2026.

\bibitem{Slepian1978}
D.~Slepian, \textit{Prolate spheroidal wave functions, Fourier
analysis, and uncertainty~-- V: The discrete case},
Bell Syst.\ Tech.\ J.\ \textbf{57} (1978), 1371--1430.

\bibitem{Osipov2013}
A.~Osipov, V.~Rokhlin, H.~Xiao,
\textit{Prolate Spheroidal Wave Functions of Order Zero},
Springer, New York, 2013.

\bibitem{Widom1964}
H.~Widom, \textit{Asymptotic behavior of the eigenvalues of
certain integral equations, II},
Arch.\ Rational Mech.\ Anal.\ \textbf{17} (1964), 215--229.

\bibitem{BerryKeating1999}
M.~V.~Berry and J.~P.~Keating, \textit{The Riemann zeros and
eigenvalue asymptotics},
SIAM Rev.\ \textbf{41} (1999), 236--266.

\bibitem{Connes1999}
A.~Connes, \textit{Trace formula in noncommutative geometry and
the zeros of the Riemann zeta function},
Selecta Math.\ (N.S.)\ \textbf{5} (1999), 29--106.

\bibitem{Weil1952}
A.~Weil, \textit{Sur les formules explicites de la th\'{e}orie
des nombres premiers},
Comm.\ Sem.\ Math.\ Univ.\ Lund, Tome Suppl.\ 1952, 252--265.

\bibitem{GoldstronMontgomery1973}
D.~A.~Goldston and H.~L.~Montgomery, \textit{Pair correlation of
zeros and primes in short intervals},
Birkh\"auser, 1987, pp.~183--203.

\bibitem{Beurling1989}
A.~Beurling, \textit{The collected works of Arne Beurling,
Vol.~2: Harmonic Analysis}, Birkh\"auser, Boston, 1989.

\bibitem{Landau1967}
H.~J.~Landau, \textit{Necessary density conditions for sampling
and interpolation of certain entire functions},
Acta Math.\ \textbf{117} (1967), 37--52.

\bibitem{PaperIV}
[Author], \textit{Spectral Phase Obstruction for Bandlimited
Operator Families: A No-Go Theorem for Zeta-Compatible Limits},
preprint, 2026.

\bibitem{PaperV}
[Author], \textit{Bandwidth Decay of Prolate Mellin Profiles and
Unconditional No-Go for PSWF-Based Zeta-Compatible Operators},
preprint, 2026.

\end{thebibliography}

\end{document}